\documentclass[11pt]{article}
\usepackage{mathblatt}
% gesetzt von werkzeuge/setzer.py; Lösungen nach bau/bauregeln.md „Lösungen“.
\newcommand{\kennung}{P9H}
\begin{document}
\pfheftstil
\fancyfoot[C]{{\footnotesize\color{mbgrau}\kennung\hfill\thepage}}
\pfheftkopf{Ist der Winkel recht?}{Klasse 9 \textperiodcentered{} Lösungen}

\begin{pfloesung}{}
\lz{1a)}{ja}{$3^2 + 4^2 = 25 = 5^2$; rechter Winkel gegenüber der längsten Seite}{}
\lz{b)}{nein}{$2^2 + 5^2 = 29 \ne 25 = 5^2$}{}
\end{pfloesung}
\begin{pfloesung}{}
\lz{2a)}{$17$ cm}{die längste Seite}{}
\lz{b)}{$1{,}1$ m}{$= 110$ cm, länger als $95$ cm}{}
\lz{c)}{$41$ cm}{$0{,}4$ m $= 40$ cm}{}
\end{pfloesung}
\begin{pfloesung}{}
\lz{3b)}{ja}{$8^2 + 15^2 = 289 = 17^2$}{}
\lz{c)}{nein}{$6^2 + 8^2 = 100 \ne 121 = 11^2$}{}
\lz{d)}{nein}{$10^2 + 11^2 = 221 \ne 225 = 15^2$; fast gleich ist nicht gleich}{}
\end{pfloesung}
\begin{pfloesung}{}
\lz{4a)}{ja, rechter Winkel bei $C$}{in cm: $7^2 + 24^2 = 625 = 25^2$; $C$ liegt gegenüber $\overline{AB}$}{}
\lz{b)}{nein}{in cm: $60^2 + 90^2 = 11\,700 \ne 12\,100 = 110^2$}{}
\end{pfloesung}
\begin{pfloesung}{}
\lz{5b)}{$24$, $45$, $51$; ja}{$24^2 + 45^2 = 2\,601 = 51^2$}{}
\lz{c)}{nein}{$3$, $4$, $6$ mal $10$: $3^2 + 4^2 = 25 \ne 36$}{}
\lz{d)}{ja}{$9^2 + 40^2 = 1\,681 = 41^2$}{}
\end{pfloesung}
\begin{pfloesung}{}
\lz{6a)}{$25$ cm}{Hypotenuse gesucht: $\sqrt{7^2 + 24^2} = \sqrt{625} = 25$}{}
\lz{b)}{$\sqrt{527} \approx 23{,}0$ cm}{Kathete gesucht: $24^2 - 7^2 = 527$ \ $\Rightarrow$ \ $\sqrt{527} \approx 22{,}96$}{}
\end{pfloesung}
\begin{pfloesung}{}
\lz{7a)}{$3{,}7$ dm}{Hypotenuse gesucht: $1{,}44 + 12{,}25 = 13{,}69$ \ $\Rightarrow$ \ $\sqrt{13{,}69} = 3{,}7$}{}
\lz{b)}{$2{,}1$ m}{Kathete gesucht: $8{,}41 - 4 = 4{,}41$ \ $\Rightarrow$ \ $\sqrt{4{,}41} = 2{,}1$}{}
\lz{c)}{$9$ cm}{Kathete gesucht, in cm: $41^2 - 40^2 = 81$ \ $\Rightarrow$ \ $\sqrt{81} = 9$}{}
\end{pfloesung}
\begin{pfloesung}{}
\lz{8.}{nein}{$1{,}2^2 + 1{,}6^2 = 4 \ne 4{,}41 = 2{,}1^2$; rechtwinklig wäre sie bei $2$ m}{}
\end{pfloesung}
\begin{pfloesung}{}
\lz{9.}{rechter Winkel bei $D$}{$\overline{AD}^2 = 32^2 + 32^2 = 2\,048 = \overline{DC}^2$; \ $2\,048 + 2\,048 = 4\,096 = 64^2 = \overline{AC}^2$ \ $\Rightarrow$ \ Umkehrung}{}
\end{pfloesung}
\end{document}
